\documentclass[10pt,a4paper]{amsart}
\usepackage[margin=19mm]{geometry}
\usepackage{amsmath,amssymb,mathtools}
\usepackage[scr=rsfs,cal=euler]{mathalfa}
\usepackage{xfrac}
\usepackage[unicode,hidelinks,bookmarks=false]{hyperref}
\newcommand{\ie}{i.e.\ }
\newcommand{\cf}{cf.\ }
\newcommand{\resp}{respectively\ }
\newcommand{\Subsection}{\subsection}
\providecommand{\nobreakdash}{\nobreak\hbox{-}}
\setlength{\emergencystretch}{2em}
\multlinegap0pt
\usepackage{mathtools}
\usepackage{xcolor}
\usepackage{algorithm}\usepackage{algpseudocode}
\usepackage{amssymb}
\newtheorem{definition}{Definition}
\newtheorem{lemma}{Lemma}
\newtheorem{theorem}{Theorem}
\newcommand{\argmin}{\mathop{\rm argmin}}
\def\x{\mathbf{x}}
\title{On the admissibility of bounds on the mean \\of discrete, scalar probability distributions\\ from an i.i.d.~sample}
\author{Erik Learned-Miller, George Bissias}
\date{Source: arXiv:2502.17223v2 (CC BY 4.0)}
\begin{document}
\maketitle

\noindent\emph{Source and licence.} On the admissibility of bounds on the mean \\of discrete, scalar probability distributions\\ from an i.i.d.~sample. Version arXiv:2502.17223v2 (CC BY 4.0), distributed by arXiv under the Creative Commons Attribution 4.0 International licence, http://creativecommons.org/licenses/by/4.0/, as recorded in the arXiv licence metadata for this exact version and shown on the abstract page https://arxiv.org/abs/2502.17223v2. Full licence notice: \texttt{licenses/arxiv-2502.17223-CC-BY-4.0.txt}.\par\smallskip

\noindent\textit{Editorial scope.} Counted unit: the article's theorem thm:conditional\_optimality (source0.tex lines 776--779). The article's complete original proof of it is delivered with it (source0.tex lines 781--817, one \texttt{proof} environment of 37 lines). No mathematics is written, repaired or re-derived: the statement and the proof are the article's own lines, in the article's own notation, and no retained line is substituted or re-flowed.  Retained uncounted as the source-local context the proof needs: lines 448--455; lines 475--478; lines 600--603; lines 722--733; lines 754--757.  Nothing was omitted from any retained span, so the package carries no omission record.  No retained line contains a citation, so the wrapper carries no bibliography and no bibliography entry.  No retained line contains a figure environment, an image-inclusion command or drawing code, so no image is delivered and the package declares no asset dependency.  Editorial formatting only, in addition: an AMS \texttt{amsart} preamble that restates the article's own declarations and macros used by the retained lines, namely the environments \texttt{definition}, \texttt{lemma}, \texttt{theorem} on separate counters, together with the standard packages those lines need.  The retained environments print in this document as Definition 1, Lemma 1, Theorem 1 in order of appearance, whereas the article prints the same items with the numbering of its own full document.  The complete native source of the article as distributed in the arXiv e-print for this exact version is bundled unchanged as \texttt{sources/173787/source0.tex}.
\begin{definition}[error set and valid set of a bound for a distribution and sample size]
\label{def:error_set_of_valid_bound}
Given a distribution $F$ with mean $\mu(F)$ over a support set $S$ and a sample size $n$, let $\Omega(S,n)$ be the sample space. Let $1-\alpha$ be the confidence level. Then the {\bf error set} of a lower bound $B$ for the distribution and the sample size is the set $E$ defined as 
$$
E(F,n) = \{\x\in \Omega(S,n): B(\x) > \mu(F)\},
$$  
that is, the subset of the sample space for which the bound is erroneous. The complement of this set, i.e., the subset for which the bound is correct, is termed the {\bf valid set}.
\end{definition}

\begin{definition}[validity of a bound for a set of distributions]
\label{def:validity_for_dists}
For a set of probability distributions $\mathcal{F}$, a confidence level $1-\alpha$ and a sample size $n$, a bound is called {\bf valid}, or alternatively, {\bf conservative}, if  it is valid for each  $F\in\mathcal{F}$. 
\end{definition}

\begin{definition}[optimality of a bound]
\label{def:opt}
    Bound $B$ is \emph{optimal} with respect to  a set of bound functions $\mathcal{B}$ if it is valid and dominates every other valid bound in $\mathcal{B}$. It is immediate that an optimal bound must be the only admissible bound in $\mathcal{B}$. 
\end{definition}

\begin{definition}[Buehler bound]
\label{def:cop}
    Consider a sample space $\Omega$, a sample order $T$, and sample $\x_{t_k} \in \Omega$, the $k$th lowest sample in order $T$. The \emph{exhasutive bound} on $\x_{t_k}$ with respect to $T$ is given by
    \begin{eqnarray}
    \label{eq:optBound}
    B^*(\x_{t_k}) =
    \begin{cases}
        \infty,& \text{if } \mathcal{F}_k=\emptyset \\
        \underset{F\in \mathcal{F}_k}{\min} \;\mu(F),              & \text{otherwise.}
    \end{cases}
    \end{eqnarray}
\end{definition}

\begin{lemma}[order-consistency of Buehler bounds]
\label{lemma:incremental_sums}
Fix total order $T$. If $\x_{t_k}$ precedes $\x_{t_j}$ in order $T$, then $B^*(\x_{t_k}) \leq B^*(\x_{t_j})$. 
\end{lemma}

\begin{theorem}[conditional optimality]
\label{thm:conditional_optimality}
Let $\mathcal{B}$ be the lower bound functions (both valid and invalid) over a sample space $\Omega$ and for a confidence level $1-\alpha$. For a given total order $T$, let $\mathcal{B}_T\subset \mathcal{B}$ be the order-consistent bounds (again, both valid and invalid) with respect to a total order $T$. Let $B^*\in \mathcal{B_T}$ be the Buehler bound of Definition~\ref{def:cop}. Then $B^*$ is optimal with respect to $\mathcal{B}_T$. 
\end{theorem}

\begin{proof}
The conditional optimality of $B^*$ requires optimality according to Definition~\ref{def:opt} among bounds in $\mathcal{B}_T$ as well as validity according to Definition~\ref{def:validity_for_dists}. We proceed by proving each separately.

{\bf Validity.} 
For any $F \in \mathcal{F}$, let $E^*(F, n)$ denote the error set for $F$ with respect to $B^*$, where by Definition~\ref{def:error_set_of_valid_bound}
\begin{eqnarray}
\label{eq:exhaustive_error_set}
    E^*(F,n) \coloneqq \{\x\in \Omega(S,n): B^*(\x) > \mu(F)\}.    
\end{eqnarray}
According to Definition~\ref{def:validity_for_dists} it will suffice to show that $\forall F \in \mathcal{F}$, $\Pr_F(E^*(F, n)) \leq \alpha$. 
To that end, suppose for the sake of contradiction that there exists $F \in \mathcal{F}$ such that $\Pr_F(E^*(F, n)) > \alpha$.
According to Lemma~\ref{lemma:incremental_sums}, $B^*$ itself will be in $\mathcal{B}_T$.
Therefore, from Equation~\ref{eq:exhaustive_error_set} we can see that there exists $\Omega_k$ such that $\Omega_k = E^*(F, n)$ and $B^*(\x_{t_k}) > \mu(F)$.
It follows then that $\Pr_F(\Omega_k) > \alpha$, which implies that $F \in \mathcal{F}_k$.
But since $\mathcal{F}_k$ is non-empty, Equation~\ref{eq:optBound} ensures that $B^*(\x_{t_k}) \leq \mu(G)$ for each $G \in \mathcal{F}_k$.
This is true in particular for $F$.
Thus, $B^*(\x_{t_k}) \leq \mu(F)$, a contradiction.
It must then be the case that $\Pr_F(E^*(F, n)) \leq \alpha$ for each $F \in \mathcal{F}$.
\color{black}
\medskip

{\bf Optimality for order $T$.}
Let $B \in \mathcal{B}_T$ be any valid bound. It suffices to show that 
\begin{eqnarray}
    B(\x_t) \leq B^*(\x_t)
\end{eqnarray}
for each $\x_t \in \Omega$. To that end, consider an arbitrary sample $\x_{t_k}$. If $B^*(\x_t) = \infty$, then the conclusion clearly holds. Otherwise, suppose that $B(\x_{t_k}) > B^*(\x_{t_k})$. It must then be the case that $\exists \delta > 0, B(\x_{t_k}) = B^*(\x_{t_k}) + \delta$.
And since $B^*$ is the Buehler bound of Definition~\ref{def:cop}, we must also have that $\forall \epsilon > 0, \exists G_\epsilon \in \mathcal{G}_k: B^*(x_{t_k}) + \epsilon \geq \mu(G_\epsilon)$.
Thus, by choosing $\epsilon < \delta$, we have that $B(\x_{t_k}) > \mu(G_\epsilon)$.
%By the definition of $B^*$, it must then be the case that $B(\x_{t_k}) > \mu(F^*)$ where $F^* = \argmin_{F \in \mathcal{F}_k} \mu(F)$.
Of course, $\Pr_{G_\epsilon}(\Omega_k) > \alpha$ since $G_\epsilon \in \mathcal{G}_k$.
And because $B$ is consistent with $T$, it must also be the case that $B(\x) > \mu(G_\epsilon)$ for all $\x \in \Omega_k$.
This, however, would imply that the error set for $G_\epsilon$ under bound $B$ has measure greater than $\alpha$, making $B$ invalid, which is a contradiction. Thus, $B(\x_t) \leq B^*(\x_t)$ for every valid $B \in \mathcal{B}_T$ and each choice of $x_t \in \Omega$. 
\color{black}

Notice that this result does not depend upon the value of the bound for any other sample, as long as the ordering is fixed. Thus, conditioning on the order makes it possible to analyze the bound one sample at a time.
\end{proof}

\end{document}
